%!TEX root = ../template.tex
%% ====================================================================
%% Capitulo 6 -- Resultados II --- Diagnostico das baselines
%% RASCUNHO COMPLETO 2026-08-30 sobre as corridas canonicas:
%%   E1+D0: outputs/runs/20260826T165010999180Z-ff9a3322b633 (D0 aceite 2026-08-26)
%%   E2:    outputs/E2/runs/20260828T153423377266Z-7926b7afbe338986221d
%%   E3:    outputs/E3/runs/20260829T122639790978Z-aa3e341622b4 (commit ad0c9db)
%%   E4:    outputs/E4/runs/20260829T135121261692Z-2ea18e7390bb (commit a016605)
%%   E6:    outputs/E6/runs/20260829T153020733605Z-ccd0312d97b3 (commit e67fdc3)
%%   D1-D9: outputs/D1-D9/runs/20260829T182000032463Z-1dde8c0c7d78 (commit 252cf7d)
%%   SENS:  outputs/SENS/runs/20260829T195138387361Z-e068552161d2 (commit 74b2247)
%%   CH6:   outputs/CH6/runs/20260829T195459405603Z-8b10ee694a83 (figuras 6.1-6.8)
%% E5 escrito como not_estimable (decisao do autor 2026-08-30); a rota E5'
%% por simulacao fica descartada para o Cap. 6.
%% Regra transversal ja aplicada: nada de aparato informatico aqui (vive no
%% Cap. 4.5); a contagem de patologias crava-se na sintese e em mais lado
%% nenhum. Figuras por incluir a partir do bundle CH6 (SVG -> PDF).
%% ====================================================================

\chapter{Resultados II --- Diagnóstico das baselines}
\label{cha:diagnostico}

Este capítulo responde à segunda questão de investigação: as baselines
energéticas em produção na unidade de destilação atmosférica são válidas para
os fins que a ISO 50006 lhes atribui --- normalizar o consumo pela carga,
sustentar a comparação de desempenho entre períodos e alimentar um sistema de
alarmes? A análise executa-se integralmente sobre o modelo de dados
validado no Capítulo~\ref{cha:plataforma}: cada resultado é uma consulta ao
mesmo esquema dimensional que alimenta o dashboard, e não a um extrato
paralelo preparado para o efeito.

A leitura do capítulo obedece à arquitetura de dois níveis definida na
Secção~\ref{sec:protocolo-baselines}. Os \emph{endpoints} confirmatórios (E1--E6) usam
inferência robusta à dependência serial --- covariância HAC e \emph{bootstrap}
de blocos --- e são os únicos de que se extraem conclusões; a bateria
descritiva (D0--D9) documenta mecanismos e contextos, com a inferência clássica
explicitamente desqualificada pela autocorrelação residual que a própria
bateria mede. Nenhuma conclusão deste capítulo assenta num $p$-value de
nível descritivo.

\section{O modelo sob análise: replicação exata}
\label{sec:e1-replicacao}

Antes de qualquer crítica, o método de produção foi replicado de forma exata a
partir das tabelas de configuração oficiais: para cada um dos cinco vectores
energéticos e cada ano civil de 2020 a 2025, a regressão $E = aQ + b$ foi
reestimada sobre os dias com carga acima do limiar oficial de
$18\,600$~t/d --- lido da configuração do modelo dimensional, nunca fixado no
código --- e comparada com os parâmetros em produção. As 30 células
vector--ano reproduzem declive, ordenada na origem, $\hat\sigma$ e $R^2$
oficiais dentro da tolerância pré-registada de $10^{-6}$, com erro máximo
observado da ordem de $10^{-10}$; a identidade dos conjuntos de dias
selecionados foi igualmente verificada célula a célula
(mecanismo e limitações na Secção~\ref{sec:reprodutibilidade}).

O ponto não é informático, é epistémico: tudo o que as secções seguintes
afirmam refere-se, com igualdade numérica verificada, ao modelo que está
efetivamente em produção --- os mesmos parâmetros que desenham as bandas de
alarme do dashboard --- e não a uma reconstrução aproximada dele.

\section{Perfil dos dados e domínio de validade}
\label{sec:d0-perfil}

%% FIGURA: outputs .../figures/D0_4_threshold_effect.svg (converter p/ PDF)
%% FIGURA: outputs .../figures/D0_6_scatter_official.svg

O painel analisado cobre $2\,403$ dias candidatos entre 2020-01-01 e o corte
oficial de 2026, com $11\,955$ observações série--dia nos cinco vectores; a
falta de dados é rara e concentrada (60 observações, em dois blocos de 30 dias
na energia elétrica, 2021-11 e 2026-07). Três características dos dados,
estabelecidas antes de qualquer inferência, condicionam a leitura de todo o
diagnóstico.

\paragraph{O limiar de carga define um domínio de validade não declarado.}
O limiar oficial de $18\,600$~t/d remove 251 dos $2\,403$ dias candidatos
(10{,}4\%) e 744 das $11\,467$ observações com energia medida. Os dias
removidos não são uma franja aleatória de dias atípicos: comparados com os
retidos, diferem de forma sistemática em todas as variáveis observadas ---
diferenças padronizadas de médias com mediana $|\mathrm{SMD}|=2{,}17$ e máximo
$10{,}5$ no consumo energético, e máximos de $19{,}7$ na temperatura de
\emph{flash} e $6{,}8$ na pressão de \emph{flash}. Trata-se de outro regime
operatório, não de uma cauda da mesma distribuição. A consequência é
estrutural: a baseline em produção só descreve a unidade em regime de carga
alta, e todo o comportamento abaixo do limiar --- arranques, reduções de
carga, operação degradada --- está fora do domínio em que o modelo alguma vez
foi estimado. Este domínio de validade implícito não está declarado em nenhum
documento do método de produção, e a sua existência é a primeira condição de
uso que o diagnóstico regista. Nos cenários de sensibilidade com o limiar a
$\pm 20\%$, a remoção varia entre 187 e 397 dias, com o mesmo padrão de
diferenças sistemáticas.

\paragraph{A energia elétrica não é uma medição diária.} Nos 77 meses
completos do período, o valor diário da energia elétrica é constante dentro de
cada mês em 77 --- um total mensal distribuído uniformemente pelos dias, não
uma medição com resolução diária. Isto explica mecanicamente o $R^2$
residual do vector ($\approx 4\times 10^{-4}$ na configuração oficial): a
regressão diária tenta explicar com a carga uma variável que, por construção,
não varia dentro do mês. Mais importante, torna o alarme diário da energia
elétrica estruturalmente cego: nenhum desvio com duração inferior a um mês é
detetável, qualquer que seja a estatística usada. Por esta razão, a energia
elétrica é excluída dos endpoints de resolução diária (E3--E5) e avaliada no
E2 apenas por um diagnóstico mensal próprio.

\paragraph{Os resíduos são fortemente dependentes.} A autocorrelação residual
de primeira ordem das regressões oficiais situa-se entre $0{,}6$ e $0{,}95$
consoante o vector e o ano. Sob esta dependência, os erros-padrão clássicos e
os $p$-values dos testes de diagnóstico habituais deixam de ser fiáveis, e a
banda de alarme $\pm 2\hat\sigma$ deixa de ter a cobertura nominal que o seu
desenho pressupõe. É este facto --- medido nos próprios dados --- que
justifica a maquinaria inferencial do protocolo
(Secção~\ref{sec:protocolo-baselines}) e que as Secções~\ref{sec:e3-estabilidade}
e~\ref{sec:e4-calibracao} traduzem em consequências operacionais.

\section{Poder preditivo contra referências triviais}
\label{sec:e2-skill}

%% FIGURAS: E2_2_delta_mae_forest.svg, E2_1_skill_heatmap.svg,
%%          E2_4_fuel_gas_forward_series.svg (bundle E2)

A primeira pergunta confirmatória é a mais elementar: os parâmetros oficiais
de um ano de treino $y$, aplicados sem reestimação ao ano seguinte $y{+}1$,
prevêem o consumo melhor do que referências triviais? A referência primária é
a média do ano de treino --- um ``modelo'' sem carga, que responde sempre o
mesmo valor. O efeito primário é a diferença emparelhada de perda absoluta
diária, $\Delta_{\mathrm{MAE}} = \mathrm{MAE}_{\mathrm{nulo}} -
\mathrm{MAE}_{\mathrm{modelo}}$, com intervalo de confiança percentil a 95\%
por \emph{moving-block bootstrap} (MBB) sobre a série das diferenças de
perda; o \emph{skill} associado é
$1 - \mathrm{MAE}_{\mathrm{modelo}}/\mathrm{MAE}_{\mathrm{nulo}}$. Não se
constrói nenhum teste agregado sobre os cinco pares anuais: cada par
vector--transição é reportado individualmente, e a conclusão é o padrão dos
intervalos. O ano incompleto de 2026 é apenas diagnóstico.

O resultado divide os vectores em dois grupos. O fuel gás tem conteúdo
preditivo real: em quatro das cinco transições completas
(2020${\to}$2021, 2022${\to}$2023, 2023${\to}$2024, 2024${\to}$2025) o
\emph{skill} situa-se entre $0{,}30$ e $0{,}40$ e o intervalo MBB de
$\Delta_{\mathrm{MAE}}$ exclui zero do lado positivo; a exceção é
2021${\to}$2022 (\emph{skill} $0{,}06$, intervalo contém zero). Conhecer a
carga vale, nas transições favoráveis, reduções de 30--40\% do erro absoluto
médio face a responder sempre a média.

Os três vectores de vapor são, para efeitos de previsão pontual, modelos
vazios. No vapor de 24~bar o \emph{skill} nunca sai do intervalo
$[-0{,}05;\,0{,}01]$; no de 10~bar, de $[-0{,}07;\,0{,}004]$; os intervalos
ora contêm zero ora o excluem por margens sem significado prático
($|\Delta_{\mathrm{MAE}}|$ até $0{,}5$~ton~GNE/d, contra consumos médios de
dezenas de toneladas). O vapor de 3~bar acrescenta instabilidade de sinal: a
transição 2020${\to}$2021 é confirmadamente \emph{pior} do que a média do
treino (\emph{skill} $-0{,}20$, intervalo exclui zero do lado negativo),
2021${\to}$2022 e 2023${\to}$2024 marginalmente melhores, as restantes
indistinguíveis de zero. Em três dos quatro vectores diários, a componente
linear em $Q$ não transporta informação utilizável de um ano para o seguinte
--- coerente com os $R^2$ oficiais destes vectores, entre $2\times10^{-5}$ e
$0{,}31$.

Duas leituras que o resultado \emph{não} autoriza. Primeiro, previsão pontual
não é normalização: um \emph{skill} nulo mata o uso da baseline como preditor,
mas não decide por si só o seu valor como normalizador de consumo pela carga
--- essa avaliação pertence às Secções~\ref{sec:e3-estabilidade}
e~\ref{sec:e4-calibracao}. Segundo, o diagnóstico de 2026 (parcial) mostra o
fuel gás com \emph{skill} $-0{,}45$ --- o modelo de 2025 é, no ano corrente,
substancialmente pior do que a média histórica --- mas o ano incompleto não
suporta conclusões; o padrão é retomado na análise de estabilidade.

A referência sazonal (mesmo dia civil do ano anterior) e a média móvel de 30
dias, referências secundárias, não alteram este quadro e constam do
Anexo~\ref{app:tabelas}. O diagnóstico mensal da energia elétrica, contra
nulos apropriados à sua resolução (média mensal do treino, mês homólogo, mês
anterior), é igualmente remetido para o anexo: com $n{=}12$ meses por ano e
alocação uniforme, nenhuma comparação atinge estatuto confirmatório.

\section{Estabilidade dos parâmetros}
\label{sec:e3-estabilidade}

%% FIGURAS: E3_1_<vector>_reference_predictions.svg (bundle E3, 4 figuras)

A recalibração anual do método de produção pressupõe que cada ano de dados
reestima \emph{a mesma} relação física entre energia e carga, a menos de ruído
amostral. O endpoint E3 testa esse pressuposto diretamente: uma regressão
empilhada 2020--2025 com interações ano$\times$carga, e um teste de Wald de
igualdade dos declives com covariância HAC (\emph{kernel} quadrático-espectral,
largura de banda de Andrews), que permanece válido sob a dependência residual
medida na Secção~\ref{sec:d0-perfil}. Os quatro testes de declive formam uma
família com correção de Benjamini--Hochberg (BH), com
Benjamini--Yekutieli (BY) como verificação conservadora; os testes de nível
formam família separada e secundária.

\begin{table}[htbp]
  \centering
  \caption[Igualdade dos declives anuais (E3)]{Teste de Wald HAC de igualdade
  dos declives anuais, 2020--2025. Família de quatro hipóteses; $q$-values BH
  e BY. Declives anuais extremos em ton~GNE/d por kton de carga.}
  \label{tab:e3-slopes}
  \begin{tabular}{lccccc}
    \toprule
    Vector & $\hat a$ mín.--máx. & $W$ & $p$ & $q_{\mathrm{BH}}$ & $q_{\mathrm{BY}}$ \\
    \midrule
    Fuel gás        & $8{,}19$ -- $11{,}75$   & $12{,}3$ & $0{,}031$            & $0{,}042$            & $0{,}087$ \\
    Vapor 24 bar    & $-0{,}15$ -- $0{,}31$   & $2{,}8$  & $0{,}731$            & $0{,}731$            & $1$ \\
    Vapor 10 bar    & $-0{,}43$ -- $0{,}15$   & $13{,}8$ & $0{,}017$            & $0{,}034$            & $0{,}071$ \\
    Vapor 3 bar     & $-0{,}15$ -- $2{,}81$   & $28{,}0$ & $3{,}7\times10^{-5}$ & $1{,}5\times10^{-4}$ & $3{,}1\times10^{-4}$ \\
    \bottomrule
  \end{tabular}
\end{table}

A igualdade dos declives é rejeitada a 5\% sob BH em três dos quatro vectores
--- vapor de 3~bar, vapor de 10~bar e fuel gás
(Tabela~\ref{tab:e3-slopes}); sob a verificação BY, robusta a qualquer
estrutura de dependência entre os testes, sobrevive a rejeição do vapor de
3~bar. Os números por trás dos testes mostram a dimensão física do problema.
No vapor de 3~bar o declive oficial percorre $+2{,}35 \to +2{,}81 \to +1{,}78
\to -0{,}15 \to +0{,}73 \to +0{,}09$: a ``mesma'' relação muda de sentido ---
em 2023, mais carga significava \emph{menos} vapor de 3~bar. No fuel gás, o
declive desce de $11{,}75$ para $8{,}19$ entre 2020 e 2024 (${-}30\%$), com a
ordenada na origem a subir de $-34$ para $+96$~ton~GNE/d em compensação. No
vapor de 10~bar o declive oscila em torno de zero entre $-0{,}43$ e $+0{,}15$.
A não-rejeição no vapor de 24~bar não é evidência de uma relação estável: com
declives anuais entre $-0{,}15$ e $+0{,}31$ em torno de zero, o que é estável
é a \emph{ausência} de relação --- o teste não distingue seis estimativas da
mesma constante de seis estimativas de nada. A família secundária dos níveis
rejeita a igualdade nos mesmos três vectores, mesmo sob BY.

A leitura operacional: a recalibração anual não está a reestimar uma constante
física com precisão crescente; está a ajustar, a cada janeiro, um modelo
diferente ao regime desse ano --- absorvendo silenciosamente no par $(a, b)$ a
deriva que um sistema de monitorização deveria expor. As predições em cargas
de referência com bandas de confiança HAC, por vector e ano, traduzem esta
instabilidade em unidades de consumo e constam das figuras deste capítulo;
a sensibilidade à largura de banda HAC (multiplicadores $0{,}5\times$ e
$2\times$) não altera nenhuma das conclusões.

\section{Calibração das bandas de alarme}
\label{sec:e4-calibracao}

%% FIGURAS: bundle E4 (cobertura por par, rajadas, forma da banda)

O sistema de alarmes em produção declara um dia anómalo quando o resíduo sai
da banda $\pm 2\hat\sigma$ do ano de treino. Se os pressupostos do desenho
valessem --- resíduos independentes, homocedásticos, aproximadamente normais,
relação estável --- a banda aplicada ao ano seguinte cobriria cerca de
$95{,}4\%$ dos dias. O endpoint E4 mede a cobertura efetiva nos cinco pares
treino${\to}$aplicação completos de cada vector diário, com intervalo de
confiança MBB a 95\%; a referência de $0{,}954$ é indicativa, e todas as
conclusões resultam da sua exclusão pelo intervalo, nunca da comparação de
pontos. O intervalo binomial clássico, que assume independência, é publicado
apenas como contraste: o intervalo MBB é $1{,}0$ a $2{,}3$ vezes mais largo em
todos os 20 pares --- a dependência serial a cobrar o seu preço à precisão.

A descalibração é generalizada e heterogénea. Em sete dos 20 pares a
referência é excluída por baixo (sub-cobertura, excesso de alarmes): fuel gás
2020${\to}$2021 (cobertura $0{,}62$), vapor de 10~bar 2020${\to}$2021,
2021${\to}$2022 e 2024${\to}$2025 ($0{,}87$, $0{,}88$, $0{,}53$) e vapor de
3~bar 2021${\to}$2022, 2022${\to}$2023 e 2024${\to}$2025 ($0{,}38$, $0{,}84$,
$0{,}64$). Em três pares é excluída por cima (sobre-cobertura, alarme
adormecido): fuel gás 2021${\to}$2022 ($0{,}98$), vapor de 10~bar e vapor de
3~bar 2023${\to}$2024 ($0{,}99$ e $0{,}997$ --- um e um alarme em 364 dias,
respetivamente). O vapor de 24~bar fica dentro do intervalo nos cinco pares,
em coerência com a Secção~\ref{sec:e3-estabilidade}: a banda de um modelo sem
relação a proteger é uma banda em torno de uma média, e essa calibra
razoavelmente.

Os casos extremos merecem tradução operacional. No vapor de 3~bar
2021${\to}$2022, cobertura $0{,}38$ significa que $62\%$ dos dias do ano
estiveram em alarme --- um sistema que grita dois dias em cada três não está a
monitorizar, está desligado na prática. No vapor de 10~bar 2024${\to}$2025,
139 dos 295 dias ($47\%$) excederam a banda; no par seguinte do mesmo vector
(2023${\to}$2024) tinham sido 2 em 364. A mesma regra de alarme, no mesmo
vector, oscila entre o silêncio quase total e o ruído permanente consoante o
par de anos --- consequência direta da instabilidade dos parâmetros e da
dependência residual que o desenho da banda ignora. O diagnóstico parcial de
2026 aponta na mesma direção, sem estatuto confirmatório: no vapor de 10~bar,
a banda de 2025 cobre $6\%$ dos dias decorridos --- 172 alarmes em 183 dias,
um colapso coerente com a subida de nível identificada no E3.

O quadro que emerge das duas últimas secções é consistente: as bandas falham
nas duas direções, e falham mais nos vectores onde os parâmetros são mais
instáveis. A análise de cobertura condicional (tercis de carga por semestre),
a distribuição das rajadas de alarme e o erro de forma da banda constante
face ao intervalo de predição correto constam do Anexo~\ref{app:tabelas}.

\section{Detectabilidade}
\label{sec:e5-detectabilidade}

O endpoint E5 perguntaria que desvio mínimo de consumo o sistema de alarmes
consegue detetar, e ao fim de quantos dias, comparando a regra
$\pm 2\hat\sigma$ em produção com alternativas sequenciais (CUSUM sobre
inovações) para tamanhos de desvio elicitados com o orientador industrial. A
pergunta é central para um sistema de monitorização --- de nada serve uma
banda calibrada que só deteta desvios de uma magnitude que nunca ocorre ---
mas não é estimável neste conjunto de dados, e o resultado publica-se como
tal, em vez de se publicar uma resposta fabricada.

A razão é estrutural. Avaliar detectabilidade exige janelas de referência
comprovadamente limpas --- períodos em que se saiba, por informação
independente das séries, que a unidade operou sem eventos --- e a marcação
dos eventos reais contra os quais medir tempos de deteção. Ambas dependem do
calendário operacional da unidade (paragens, arranques, campanhas,
intervenções), pedido em 2026-08-18 e confirmado indisponível em 2026-08-25.
Inferir as paragens a partir das próprias séries de consumo, como atalho,
contaminaria circularmente o exercício: os ``eventos'' seriam definidos pela
mesma estatística cuja capacidade de deteção se pretende medir. Pela mesma
causa raiz ficam não estimáveis o diagnóstico descritivo de eventos
operacionais (D4) e a componente pós-arranque da análise de suporte (D6). A
indisponibilidade fica declarada como limitação
(Secção~\ref{sec:limitacoes}); o desenho completo do endpoint permanece
especificado no protocolo, executável no dia em que o calendário exista.

O que as secções anteriores já mostram, sem substituir o E5: um sistema cuja
cobertura efetiva oscila entre $0{,}38$ e $0{,}997$ consoante o vector e o par
de anos (Secção~\ref{sec:e4-calibracao}) não tem uma curva de deteção única
--- tem uma por regime de descalibração, e nenhuma delas é a nominal.


\section{Suficiência da variável explicativa}
\label{sec:e6-suficiencia}

%% FIGURAS: bundle E6 (E6_1_<vector>_skill.svg, 4 figuras)

Se a carga sozinha explica pouco (Secção~\ref{sec:e2-skill}), acrescentar
variáveis de processo resolve? O endpoint E6 responde com a única forma de
resposta que não se auto-engana: valor preditivo \emph{incremental},
estritamente fora da amostra. Para cada vector diário e ano de treino,
estimam-se dois modelos encaixados --- $E = aQ + b$ e $E = aQ + \gamma'Z + b$,
com $Z$ pré-especificado antes de qualquer execução (densidade da carga,
temperatura e pressão de \emph{flash}; o grau API é excluído por ser
transformação determinista da densidade) --- e comparam-se as perdas absolutas
no ano seguinte, sem reestimação. O efeito é a diferença emparelhada
$\Delta_{\mathrm{MAE}} = \mathrm{MAE}_{Q} - \mathrm{MAE}_{Q+Z}$ com intervalo
MBB a 95\%: positivo com exclusão de zero, as covariáveis acrescentam;
negativo, prejudicam. O ganho \emph{in-sample}, que é garantido por
construção, não entra no critério.

\begin{table}[htbp]
  \centering
  \caption[Valor incremental das covariáveis (E6)]{Valor preditivo
  incremental das covariáveis pré-especificadas:
  $\Delta_{\mathrm{MAE}} = \mathrm{MAE}_{Q} - \mathrm{MAE}_{Q+Z}$ em
  ton~GNE/d, por transição treino${\to}$aplicação. Marcas: $^{+}$ intervalo
  MBB a 95\% exclui zero pelo lado positivo (covariáveis confirmadamente
  úteis); $^{-}$ exclui zero pelo lado negativo (confirmadamente
  prejudiciais); sem marca, o intervalo contém zero.}
  \label{tab:e6-incremental}
  \begin{tabular}{lrrrrr}
    \toprule
    Vector & \multicolumn{1}{c}{20${\to}$21} & \multicolumn{1}{c}{21${\to}$22}
    & \multicolumn{1}{c}{22${\to}$23} & \multicolumn{1}{c}{23${\to}$24}
    & \multicolumn{1}{c}{24${\to}$25} \\
    \midrule
    Fuel gás     & $4{,}16^{+}$  & $-3{,}65^{-}$ & $-1{,}05$      & $-0{,}66$     & $-7{,}09^{-}$ \\
    Vapor 24 bar & $-0{,}18^{-}$ & $-0{,}33^{-}$ & $-1{,}09^{-}$  & $-0{,}60^{-}$ & $0{,}26^{+}$  \\
    Vapor 10 bar & $0{,}09^{+}$  & $1{,}31^{+}$  & $0{,}17$       & $0{,}32^{+}$  & $0{,}36^{+}$  \\
    Vapor 3 bar  & $-0{,}01$     & $0{,}05$      & $-4{,}18^{-}$  & $8{,}26^{+}$  & $0{,}83^{+}$  \\
    \bottomrule
  \end{tabular}
\end{table}

O resultado (Tabela~\ref{tab:e6-incremental}) desautoriza as duas respostas
simples. As covariáveis não são inúteis: no vapor de 10~bar acrescentam valor
confirmado em quatro das cinco transições, com reduções de erro entre $2\%$ e
$18\%$ e nenhum par negativo --- o único caso em que o degrau multivariado
que a ISO 50006 sanciona teria ajudado de forma consistente. Mas também não
são um remédio: no fuel gás, o vector com melhor modelo em $Q$, as
covariáveis \emph{pioram} a previsão em duas transições com exclusão de zero
--- na transição 2024${\to}$2025 o erro absoluto médio agrava-se $32\%$ --- e
no vapor de 24~bar prejudicam em quatro das cinco. O vapor de 3~bar exibe o
padrão mais instrutivo: $+42\%$ de redução de erro em 2023${\to}$2024,
prejuízo confirmado em 2022${\to}$2023, nada nas restantes --- o mesmo
conjunto de covariáveis oscila entre valioso e nocivo consoante o par de
anos.

O mecanismo é o mesmo diagnosticado na Secção~\ref{sec:e3-estabilidade},
agora na dimensão das covariáveis: coeficientes $\gamma$ ajustados ao regime
do ano de treino não sobrevivem à mudança de regime do ano seguinte ---
\emph{overfitting} para a frente. Onde a relação com as covariáveis é
fisicamente estável (10~bar), o valor incremental existe e repete-se; onde
não é, as covariáveis importam para dentro do modelo mais uma fonte de
não-estacionariedade. A consequência prescritiva, desenvolvida no
Capítulo~\ref{cha:discussao}: a admissão de covariáveis numa baseline não
pode ser decisão de catálogo --- exige demonstração forward caso a caso, e a
análise de sensibilidade da Secção~\ref{sec:sintese-modos-falha} mostra que
mesmo o caso favorável do 10~bar é menos sólido do que esta tabela sugere. A
comparação secundária de cobertura das bandas dos dois modelos e os
coeficientes estimados constam do Anexo~\ref{app:tabelas}.


\section{Diagnósticos complementares}
\label{sec:diagnosticos-desc}

%% FIGURAS: bundle D1-D9 (8 figuras compostas; seleccionar D1, D5, D6, D9)

A bateria descritiva D1--D9 fecha o diagnóstico documentando os mecanismos
por trás dos endpoints. O seu estatuto é deliberadamente inferior: todos os
$p$-values desta secção são índices descritivos calculados sob pressupostos
de independência que a própria bateria mostra serem falsos, e nenhum rótulo
aqui é conclusão. As tabelas completas constam do Anexo~\ref{app:tabelas};
esta secção destaca os quatro resultados que sustentam ou contextualizam as
secções anteriores.

\paragraph{A dependência residual é universal e forte (D1--D2).} Medida
sistematicamente sobre os resíduos das regressões oficiais, a autocorrelação
de primeira ordem em desfasamento civil situa-se entre $0{,}45$ e $0{,}90$
nos 24 anos-vector completos (estatística de Durbin--Watson entre $0{,}15$ e
$1{,}08$), com os valores mais altos e menos variáveis no vapor de 3~bar
($0{,}76$--$0{,}90$) --- confirmando, vector a vector e ano a ano, o perfil
preliminar da Secção~\ref{sec:d0-perfil}. A estrutura de variância acompanha:
os índices descritivos de heterocedasticidade (Breusch--Pagan, ARCH em
desfasamentos civis) são elevados na generalidade das células, e o desvio
móvel de 30 dias oscila, dentro do mesmo ano, por fatores que a banda
constante $\pm 2\hat\sigma$ não acomoda. A sazonalidade mensal existe mas é
modesta em quota de variância ($\eta^2$ entre $0{,}03$ e $0{,}09$) ---
estrutura temporal omitida, não explicação dominante. Este bloco é o
fundamento empírico de toda a maquinaria HAC/MBB do capítulo: não foi uma
preferência metodológica, foi uma imposição dos dados.

\paragraph{Não há regimes operatórios discretos a que atribuir a
instabilidade (D5).} No espaço operacional padronizado (carga e covariáveis),
o melhor agrupamento $k$-\emph{means} é $k{=}2$ com \emph{silhouette}
$0{,}35$ --- fraco --- e o DBSCAN pré-parametrizado devolve um único
\emph{cluster} com 71 dias de ruído. O espaço operacional dos anos completos
é essencialmente unimodal: a instabilidade dos parâmetros
(Secção~\ref{sec:e3-estabilidade}) não é o artefacto de dois regimes óbvios
que uma segmentação simples separaria, e uma família de baselines por regime
não é identificável a partir dos dados sem o calendário operacional.

\paragraph{A avaliação forward é interpolação, não extrapolação (D6).} A
sobreposição das densidades de carga entre anos consecutivos
(índice de \emph{overlap} de núcleos, 0 a 1) situa-se entre $0{,}79$ e
$0{,}83$ nos cinco pares completos, com $0{,}79$ no par diagnóstico
2025${\to}$2026. Os anos de aplicação vivem, em grande medida, dentro do
suporte de carga dos anos de treino: as falhas de estabilidade e de
calibração das Secções~\ref{sec:e3-estabilidade}
e~\ref{sec:e4-calibracao} ocorrem \emph{no interior} do domínio amostrado, e
não podem ser desculpadas como extrapolação para cargas nunca vistas.

\paragraph{Os vectores não falham de forma independente (D9).} As correlações
cruzadas dos resíduos padronizados entre vectores são aproximadamente nulas
em todos os pares excepto um: vapor de 24~bar contra vapor de 10~bar, com
$r = -0{,}58$ agregado nos anos completos, entre $-0{,}36$ e $-0{,}82$
consoante o ano, e $-0{,}93$ no parcial de 2026 --- uma anti-correlação forte
e persistente, compatível com o despacho entre níveis de pressão da rede de
vapor: o que um nível consome a mais, o outro tende a consumir a menos.
Registam-se ainda 77 dias com desvios simultâneos $|z| > 2$ em dois ou mais
vectores. A consequência para o desenho de alarmes é directa: alarmes
simultâneos em vectores acoplados não são evidência independente, e qualquer
agregação de alarmes que os some como se o fossem sobrestima a informação
disponível.

Os restantes blocos confirmam sem alterar o quadro: as caudas dos resíduos
são moderadamente assimétricas (assimetria entre $-2{,}1$ e $+2{,}2$, com
intervalos MBB largos), nenhuma semana ISO domina qualquer ajuste anual
(efeito máximo da remoção de uma semana: $0{,}27\hat\sigma$ na predição, em
1\,284 remoções), e a verificação de plausibilidade física regista 8 declives
anuais negativos em 35 células e 10 trocas de sinal entre anos completos
consecutivos --- fisicamente implausíveis para relações energia--carga e
coerentes com a identificação fraca diagnosticada em E2/E3 --- com as gamas
de engenharia formais declaradas indisponíveis. O diagnóstico de eventos
operacionais (D4) e a análise pós-arranque ficam não estimáveis pela
indisponibilidade do calendário (Secção~\ref{sec:e5-detectabilidade}).


\section{Síntese: os modos de falha}
\label{sec:sintese-modos-falha}

%% FIGURA: SENS_survival_matrix.svg (bundle SENS, corrida canonica T10)

Nenhum vector passa o diagnóstico completo --- mas nenhum falha da mesma
maneira, e é a diversidade dos modos de falha, mais do que qualquer falha
individual, que condena a receita única em produção. A
Tabela~\ref{tab:sintese-vectores} cruza os quatro vectores diários com as
quatro dimensões confirmatórias; a leitura por linhas identifica cinco
patologias qualitativamente distintas.

\begin{table}[htbp]
  \centering
  \caption[Síntese do diagnóstico por vector]{Síntese do diagnóstico:
  conclusões confirmatórias por vector e dimensão. E2: \emph{skill} contra a
  média do treino; E3: igualdade dos declives anuais (BH a 5\%); E4:
  cobertura da banda $\pm 2\hat\sigma$ nos cinco pares; E6: valor incremental
  forward das covariáveis.}
  \label{tab:sintese-vectores}
  \small
  \begin{tabular}{lp{2.5cm}p{2.6cm}p{2.9cm}p{2.9cm}}
    \toprule
    Vector & E2 --- skill & E3 --- estabilidade & E4 --- calibração & E6 --- covariáveis \\
    \midrule
    Fuel gás & real (0{,}30--0{,}40 em 4/5) & rejeitada & falha em 2 pares (0{,}62; 0{,}98) & prejudicam (até $-32\%$) \\
    Vapor 24 bar & nulo; MM30 domina & não rejeitada & dentro do IC (5/5) & prejudicam (4/5) \\
    Vapor 10 bar & nulo & rejeitada; nível sobe em 2026 & falha em 4 pares; colapso YTD (0{,}06) & úteis (4/5), mas frágeis \\
    Vapor 3 bar & errático (sinal instável) & rejeitada (sobrevive a BY); declive muda de sinal & a pior (0{,}38--0{,}997, duas direções) & erráticas ($-14\%$ a $+42\%$) \\
    \bottomrule
  \end{tabular}
\end{table}

\paragraph{As cinco patologias.} O \emph{fuel gás} é o único vector com
modelo genuíno --- e é precisamente por ter uma relação real que a sua
patologia é visível: deriva paramétrica (declive $-30\%$ em cinco anos)
absorvida pela recalibração anual, com bandas que falham nos pares que
atravessam a deriva. O \emph{vapor de 24~bar} é o caso insidioso: passa nos
testes de estabilidade e de calibração porque não há relação nenhuma a
falhar --- declives anuais em torno de zero, \emph{skill} nulo, uma média
disfarçada de modelo. É saúde aparente por vacuidade, e é o vector que mais
enganaria uma auditoria superficial. O \emph{vapor de 10~bar} tem uma
mudança de regime em curso: nível a subir, banda de 2025 a cobrir $6\%$ dos
dias de 2026 --- o modelo já não descreve a unidade que existe. O
\emph{vapor de 3~bar} falha em todas as dimensões ao mesmo tempo, com o
declive a trocar de sinal entre anos --- identificação fraca sobre uma
relação que provavelmente nunca foi aproximadamente linear e estática. A
\emph{energia elétrica} falha antes de qualquer estatística: um alarme
diário sobre uma alocação mensal uniforme é estruturalmente cego por
construção (Secção~\ref{sec:d0-perfil}).

\paragraph{O mecanismo comum: a recalibração anual mascara aquilo que devia
expor.} As Secções~\ref{sec:e3-estabilidade} e~\ref{sec:e4-calibracao} são
duas faces do mesmo fenómeno. Cada janeiro, o refit anual reabsorve no par
$(a,b)$ a deriva acumulada no ano anterior --- e o sistema de alarmes,
recalibrado sobre o regime já derivado, volta a tratar como normalidade o
que um sistema de monitorização deveria assinalar. A consequência é
perversa: quanto mais lento e estrutural o fenómeno (degradação, incrustação,
desativação), mais eficazmente a recalibração o esconde; sobram para os
alarmes os desvios rápidos, precisamente os que a dependência residual e a
descalibração das bandas tornam pouco fiáveis. O sistema é assim cego por
construção aos fenómenos lentos que mais importam e ruidoso nos rápidos que
menos importam.

\paragraph{Sensibilidade: o que sobrevive a mexer nas convenções.} Todo o
diagnóstico assenta em duas convenções herdadas do método de produção: o
limiar de carga e a janela anual civil. Uma grelha de sensibilidade
pré-registada reavaliou as 104 conclusões confirmatórias
(E2$\times$nulos, E3, E4, E6) sob seis combinações --- limiar a $0{,}8\times$,
$1{,}0\times$ e $1{,}2\times$, janelas civis e janelas de 12 meses ancoradas
a 1 de julho --- num total de 624 reavaliações, com a célula de referência a
reproduzir exatamente os testes deterministas do capítulo. Ao limiar, o
diagnóstico é moderadamente robusto: nas janelas civis, apenas 21 das 312
reavaliações mudam de rótulo, quase todas perdas de significância em pares
fronteiriços; as exceções substantivas ficam identificadas --- a rejeição de
estabilidade do fuel gás não sobrevive ao limiar reduzido, a descalibração
do vapor de 3~bar em 2023${\to}$2024 degenera (a sub-cobertura vive nos dias
junto ao limiar), e o valor incremental do vapor de 10~bar
(Secção~\ref{sec:e6-suficiencia}) cai em três células e num
re-sorteio da própria célula de referência --- é a conclusão mais frágil do
capítulo, e fica qualificada como tal.

A janela é outra história: re-ancorar os doze meses de janeiro para julho
muda 133 das 312 reavaliações ($43\%$). O caso mais informativo é a própria
estabilidade: com janelas julho--junho, o vapor de 24~bar --- estável em
janelas civis --- passa a rejeitar a igualdade dos declives nas três células,
e o vapor de 10~bar deixa de a rejeitar, com a subida de nível de 2026
diluída entre duas janelas. O veredicto de estabilidade de uma baseline
anual depende, em parte mensurável, de uma convenção administrativa --- o mês
em que a janela começa --- e não apenas da física do processo. Para um
método que a ISO 50006 trata como neutro, é um resultado central: a escolha
da janela é uma decisão de modelação com consequências, e as guidelines do
Capítulo~\ref{cha:discussao} tratam-na como tal.

\paragraph{Resposta à questão de investigação.} As baselines em produção na
unidade não são válidas para os fins que a ISO 50006 lhes atribui. Não por
um defeito único e corrigível, mas por cinco modos de falha distintos que
partilham uma causa: um modelo linear estático, reestimado anualmente sobre
janelas convencionais, aplicado a um processo contínuo não-estacionário cujo
comportamento os dados --- agora reconciliados e auditáveis --- documentam
com precisão. O Capítulo~\ref{cha:discussao} converte cada patologia numa
prescrição; o que este capítulo estabelece é que nenhuma prescrição pode
assumir o que estas baselines assumiram.

